\section{Conclusion}
\label{sec:conclusion}

EPSS-only vulnerability prioritization decays from mean P@50 $= 0.331$ at Window~1 to $0.034$ at Window~6 --- an 89.7\% relative drop --- under the simulated bi-weekly patching cadence on the EEHDA synthetic fleet. HygienePrio-full, which integrates host-level hygiene signals with EPSS, maintains mean P@50 $\geq 0.499$ at every window and outperforms EPSS-only in all 150 of 150 window-seed pairs. By Window~2, even the simplest hygiene baseline (HRS-only at $0.229$) outperforms exploit-likelihood scoring ($0.149$). The temporal finding strengthens the case for hygiene augmentation beyond what single-window evaluations show: hygiene signals are not a supplement to EPSS but become its replacement as the high-EPSS CVE backlog is exhausted.

The practical implication is concrete. Organisations that rely on EPSS as their primary prioritization signal should expect degrading returns after roughly four to six weeks of active patching at typical capacities and should supplement EPSS with host-hygiene signals to maintain scheduling quality over time. This implication is bounded to the synthetic evaluation context; external validation on real fleet telemetry with longitudinal exploitation outcome data --- not the HRS-derived proxy used here --- is required before any deployment recommendation.

The temporal lens also reframes prior single-window findings. The null result for CVE-level augmentation in~\cite{paper1vulnprio} was measured in the regime where EPSS is at its strongest; the positive result for hygiene augmentation in~\cite{paper4hygieneprio} was likewise measured at the window most favourable to EPSS. The present paper shows that the operational case for hygiene-augmented scoring is strongest precisely in the regime that single-window evaluations cannot reach.

\textbf{Reproducibility.} The simulator, scoring code, calibration script, pre-registration protocol, and frozen results are available at \url{https://github.com/Harshavardhanmalla-Labs/research-papers/tree/main/paper5}.

A secondary finding is that the pre-registered hypothesis of fixed-weight sufficiency (H3) was rejected: per-window weight recalibration on the five calibration seeds adds $+17$ to $+21$ pp at Windows~1--2 and $+9$ to $+10$ pp at Windows~3--4 over the HygienePrio fixed weights. Operations teams running HygienePrio-style scoring should refit the scorer on rolling calibration history rather than freeze at Window~1.

\textbf{Future work.} The most important next step is evaluation on real, appropriately anonymised fleet telemetry using exploitation event logs (not HRS) as ground truth. Secondary directions include: (i) characterising the EPSS decay rate as a function of patching capacity ($K$) and new-CVE arrival rate, (ii) replacing the coarse grid recalibration with online weight learning that respects the pre-registration discipline (e.g.\ rolling cross-validation with a frozen objective), and (iii) evaluating non-greedy multi-window optimisation against the greedy single-window scorer studied here.
